\documentclass[11pt,a4paper]{article}
\usepackage{tex/csmi}
\usepackage{hyperref}

\usepackage{listings,lstautogobble}

\usepackage[many]{tcolorbox}
\usepackage{xcolor}
%\usepackage[skins,listings,breakable]{tcolorbox}
\tcbuselibrary{listings}

\usepackage{accsupp}
%\lstset{language=c++,showspaces=false,showstringspaces=false,breaklines=true,columns=fullflexible,keepspaces=true,captionpos=t,literate={>>}{\ensuremath{>>}}1,mathescape,
%  numbers=left,numberstyle=\color{black},stepnumber=1,tabsize=1,numbersep=-5pt,framexleftmargin=-10pt,%xleftmargin=5ex
%}

\definecolor{vertfonce}{rgb}{0.,0.5,0.}
\definecolor{cppcommentcolor}{rgb}{0.,0.5,0.}
\definecolor{cppstringcolor}{rgb}{0.6,0.1,0.1}

\newcommand{\noncopynumber}[1]{%
    \BeginAccSupp{method=escape,ActualText={}}%
    #1%
    \EndAccSupp{}%
}

%\lstset{breakindent=10pt,autobreakindent=0pt}
\lstdefinestyle{mycppstyle}{
  language=C++,
  showspaces=false,showstringspaces=false,breaklines=true,columns=fullflexible,keepspaces=true,%captionpos=t,
  literate={>>}{\ensuremath{>>}}1,mathescape,
  numbers=left,numberstyle=\color{black},stepnumber=1,numbersep=5pt,xleftmargin=1ex,
  numberstyle=\noncopynumber,
  frame=none,%single,%leftline,
  %breaklines=true,
  %columns=fullflexible,
  basicstyle=\footnotesize\ttfamily,% \footnotesize\bf\ttfamily,
  %autogobble=true,
  %gobble=8,
  stringstyle=\color{cppstringcolor},%
  keywordstyle=\color{CSMITeal},
  commentstyle=\ttfamily\color{cppcommentcolor},
  literate={é}{{\'e}}1  {è}{{\`e}}1 {à}{{\`a}}1
  %inputencoding=latin1,
  %extendedchars=true,% permet d'avoir des accents dans le code
  %caption={C++ code using listing}
  %tabsize=2
}

\lstdefinestyle{mycppstyletiny}{
  language=C++,
  showspaces=false,showstringspaces=false,breaklines=true,columns=fullflexible,keepspaces=true,%captionpos=t,
  literate={>>}{\ensuremath{>>}}1,mathescape,
  numbers=left,numberstyle=\color{black},stepnumber=1,numbersep=5pt,xleftmargin=1ex,
  numberstyle=\noncopynumber,
  frame=none,%single,%leftline,
  basicstyle=\scriptsize\ttfamily,% \footnotesize\bf\ttfamily,
  %autogobble=true,
  %gobble=8,
  stringstyle=\color{cppstringcolor},%
  keywordstyle=\color{CSMITeal},
  commentstyle=\ttfamily\color{cppcommentcolor},
  literate={é}{{\'e}}1  {è}{{\`e}}1 {à}{{\`a}}1
  %inputencoding=latin1,
  %extendedchars=true,% permet d'avoir des accents dans le code
  %caption={C++ code using listing}
  %tabsize=2
}

\newtcblisting{mycpplisting}[2][]{
    arc=0pt, outer arc=0pt,
    listing only,
    colback=CSMILight,
    colbacktitle=CSMINavy,top=0mm,bottom=0mm,boxsep=0mm,middle=0mm,boxrule=0.1pt,
    enhanced,breakable,%left=15.5pt,
    listing remove caption=false,
    overlay={
      \fill[gray!30]
      ([xshift=-0pt]frame.south west)
      rectangle
      ([xshift=13.5pt]frame.north west);
    },
    listing style=mycppstyle,
    title=#2,
    #1
}
\newtcbinputlisting{\mycpplistingFromFile}[3][]{
    arc=0pt, outer arc=0pt,
      listing file={#2},
    listing only,
    colback=CSMILight,
    colbacktitle=CSMINavy,top=0mm,bottom=0mm,boxsep=0mm,middle=0mm,boxrule=0.1pt,
    enhanced,breakable,%left=15.5pt,
    listing remove caption=false,
    overlay={
      \fill[gray!30]
      ([xshift=-0pt]frame.south west)
      rectangle
      ([xshift=13.5pt]frame.north west);
    },
    listing style=mycppstyle,
    title=#3,
    #1
}



\lstdefinestyle{mysyntaxstyle}{
  basicstyle=\footnotesize\ttfamily
}


\newcommand{\mycpptext}[1]{{\ttfamily\bf{#1}}} %  {\ensuremath{\mathbb{P}_{#1}\mathbb{P}_{#2}\mathbb{G}_{#3}}\xspace}



\definecolor{shellpromptcolor}{rgb}{0.,0.5,1}

\newtcblisting{commandshell}{
  colback=black,colupper=white,colframe=CSMITeal,
  listing only,%size=fbox,%minimal,
  listing options={
    style=tcblatex,
    basicstyle=\linespread{1.1}\normalfont\ttfamily\footnotesize\color{white},
    language=sh,backgroundcolor=\color{black},
    %numbers=left,numberstyle=\color{white},
    %escapeinside     = {/*@}{@*/}, % Allow LaTeX inside these special comments
    %  mathescape,
    escapechar={@+},
    literate={\$prompt\$}{{\normalfont\footnotesize\textcolor{shellpromptcolor}{user \$}}}1 {é}{{\'e}}1  {è}{{\`e}}1 {à}{{\`a}}1
  },
  %every listing line={\textcolor{red}{\small\ttfamily\bfseries \$> }}
}

%%%%%%%%%%%%%%%%%%%%%%%%%%%%%%%%%%%%%%%%%%%%%%%%%%%%%%%%%%%%%%%%%%%%%%%%%%%%%%%%%%%%%%%%%%%%%%%%%
% algorithmic
%%%%%%%%%%%%%%%%%%%%%%%%%%%%%%%%%%%%%%%%%%%%%%%%%%%%%%%%%%%%%%%%%%%%%%%%%%%%%%%%%%%%%%%%%%%%%%%%%
%\usepackage[plain]{algorithm}
\usepackage{algorithm}
\usepackage{algpseudocode}

\makeatletter
\renewcommand{\ALG@beginalgorithmic}{\small}
\makeatother
\algrenewcommand\algorithmicindent{1.0em}%
%\algrenewcommand\algorithmicindent{0.5em}
\algnewcommand\And{\textbf{ and }}
\algnewcommand\Or{\textbf{ or }}
\algnewcommand\To{\textbf{ to }}

% begin vertical rule patch for algorithmicx (http://tex.stackexchange.com/questions/144840/vertical-loop-block-lines-in-algorithmicx-with-noend-option)
\makeatletter

% start with some helper code
% This is the vertical rule that is inserted
    \newcommand*{\algrule}[1][\algorithmicindent]{\makebox[#1][l]{\hspace*{.5em}\thealgruleextra\vrule height \thealgruleheight depth \thealgruledepth}}%
% its height and depth need to be adjustable
\newcommand*{\thealgruleextra}{}
\newcommand*{\thealgruleheight}{.75\baselineskip}
\newcommand*{\thealgruledepth}{.25\baselineskip}
\newcount\ALG@printindent@tempcnta
\def\ALG@printindent{%
    \ifnum \theALG@nested>0% is there anything to print
        \ifx\ALG@text\ALG@x@notext% is this an end group without any text?
            % do nothing
        \else
            \unskip
            \addvspace{-1pt}% FUDGE to make the rules line up
            % draw a rule for each indent level
            \ALG@printindent@tempcnta=1
            \loop
                \algrule[\csname ALG@ind@\the\ALG@printindent@tempcnta\endcsname]%
                \advance \ALG@printindent@tempcnta 1
            \ifnum \ALG@printindent@tempcnta<\numexpr\theALG@nested+1\relax% can't do <=, so add one to RHS and use < instead
            \repeat
        \fi
    \fi
}%
%\usepackage{etoolbox}
% the following line injects our new indent handling code in place of the default spacing
\patchcmd{\ALG@doentity}{\noindent\hskip\ALG@tlm}{\ALG@printindent}{}{\errmessage{failed to patch}}
\makeatother

% New definitions
\algnewcommand\algorithmicswitch{\textbf{switch}}
\algnewcommand\algorithmicendswitch{\textbf{switch}}
\algnewcommand\algorithmiccase{\textbf{case}}
\algnewcommand\algorithmicassert{\texttt{assert}}
\algnewcommand\Assert[1]{\State \algorithmicassert(#1)}%
\algdef{SE}[SWITCH]{Switch}{EndSwitch}[1]{\algorithmicswitch\ #1\ \algorithmicdo}{\algorithmicend\ \algorithmicswitch}%
\algdef{SE}[CASE]{Case}{EndCase}[1]{\algorithmiccase\ #1}{\algorithmicend\ \algorithmiccase}%
%\algtext*{EndSwitch}%
\algtext*{EndCase}%

\lstdefinestyle{mycppstyleInTabular}{
  style=mycppstyle,numbers=none,basicstyle=\footnotesize\ttfamily,
  xleftmargin=0pt,mathescape=true
}
\newtcblisting{mycpplistingInTabular}[2][]{
  enhanced,sharp corners,listing only,
  colback=CSMILight,colframe=CSMIRule,boxrule=0.4pt,
  colbacktitle=CSMINavy,coltitle=white,fonttitle=\small\sffamily\bfseries,
  left=7pt,right=7pt,top=6pt,bottom=6pt,
  listing style=mycppstyleInTabular,title={#2},#1
}

\renewcommand{\paragraph}[1]{%
  \par\addvspace{2.5ex plus 1ex minus .2ex}%
  {\normalfont\normalsize\bfseries #1}\par
  \addvspace{1.3ex plus .2ex}%
}
\renewcommand{\CSMIdoctype}{RAPPORT DU DEVOIR MAISON}
\renewcommand{\CSMIshorttitle}{Sujet 3 / Rapport}
\renewcommand{\CSMIsubtitle}{Calcul scientifique 2}
\renewcommand{\CSMIcovernote}{ Pourquoi la cuillère en métal devient-elle chaude au manche dans la soupe, alors que la cuillère en bois ne le devient pas ? Étude sous l'angle de la modélisation et de la résolution numérique.}
\newcommand{\dollar}{\mbox{\textdollar}}
\renewcommand{\familydefault}{\sfdefault}

\title{Diffusion de la chaleur dans le manche d'une cuillère}
\author{Pauline Wierzbicki / M1 CSMI}
\date{\today}

\begin{document}

\maketitle

\tableofcontents
\clearpage

\section{Introduction}
\label{sec:introduction}
% Présentation du sujet, de la problématique et des objectifs.
\subsection{Mise en situation}
Ce devoir maison trouve sa motivation dans une observation quotidienne : 
lorsque l'on plonge une cuillère dans une soupe chaude, sa température varie selon le matériau qui la compose.
Avec une cuillère en métal, la chaleur de la soupe se propage rapidement le long du manche, le rendant brûlant au toucher.
À l'inverse, une cuillère en bois reste tempérée à l'extrémité.

Dès lors, comment expliquer physiquement cette différence de comportement ?
C'est tout l'objet de ce travail, qui vise à modéliser et quantifier ce transfert thermique par une implémentation en Python.

\subsection{Objectifs}
Les objectifs de ce devoir maison sont :

\begin{itemize}
  \item Modéliser le problème de la cuillère par l'équation de la chaleur en une dimension d'espace et de temps
  \item Implémenter cette modélisation en Python, en utilisant un stockage creux de matrices
  \item Utiliser plusieurs méthodes de résolution (directes, itératives) et vérifier la convergence des méthodes itératives
  \item Afficher des résultats sous forme graphique et les expliquer
  \item Comparer les résultats obtenus en modélisant une cuillère en métal, et une cuillère en bois
\end{itemize}

\section{Formulation du problème}
\label{sec:formulation_mathematique}
% Équation de la chaleur 1D, paramètres physiques (diffusivité du métal et du bois), conditions initiales et aux limites.
Dans cette section nous allons rapidement donner la forme de l'équation de la chaleur ainsi que sa discrétisation.
Ces résultats sont issus du cours de Calcul Scientifique 1 du M1S1 CSMI.

L'équation de la chaleur en une dimension d'espace et une dimension de temps est de la forme :

\begin{equation}
  \frac{\partial u}{\partial t} (x,t) = \alpha \frac{\partial^2 u}{\partial x^2} (x,t)
  \label{eq:diffusion}
\end{equation}
Avec $ x \in [a,b] \subset \mathbb{R} $ et $ t \in [0,T] \subset R^+ $

Le choix de la 1D en espace s'explique par le fait que le manche d'une cuillère peut être apparenté à une barre,
et donc à un problème en une dimension d'espace.

Pour compléter l'équation (\ref{eq:diffusion}), nous devons spécifier une condition initale ainsi que des conditions aux limites.
Nous utiliserons une condition initiale constante, en partant du principe que le manche de la cuillère, 
au moment d'être immergé dans la soupe, est à température ambiante (disons 21°C).

Pour ce qui est des conditions aux limites, nous utiliserons celles de Dirichlet, qui consistent à fixer une température
constante à chaque bords. 
Le bout du manche immergé aura une température égale à celle de la soupe (disons 80°C), 
et le bout exposé à l'air aura la même température que la condition initiale tout du long (21°C pour la température ambiante).

Le choix de conditions aux limites de Dirichlet est un peu simpliste.
En effet, elles ne représentent pas completement la réalité, car les vapeurs chaudes émanées par la soupe chauffent indirectement le manche.
Nous ne nous intéresserons uniquement au phénomène de diffusion dans le manche par le contact d'une source chaude (ici la soupe), qui 
ne refroidit pas dans le temps.

Finalement, nous nous retrouvons avec ce problème :

\begin{equation}
\label{eq:chaleur-generale}
\begin{cases}
  \frac{\partial u}{\partial t} (x,t) = \alpha \frac{\partial^2 u}{\partial x^2} (x,t), & x \in [a,b],\; t \in [0,T] \\
  u(x,0) = u_{\text{init}}(x), & x \in [a,b] \\
  u(a,t) = T_{\text{soupe}}, & t \in [0,T] \\
  u(b,t) = T_{\text{ambiante}}, & t \in [0,T]
\end{cases}
\end{equation}

Les valeurs de $u_{\text{init}}(x)$, $T_{\text{soupe}}$ et $T_{\text{ambiante}}$ seront spécifiées lors de la modélisation.

\subsection{Discrétisation}

En se servant du cours de Calcul scientifique 1, nous trouvons deux discrétisations sous forme matricielle :

\begin{equation}
\label{discretisation:explicite}
  U^{n+1} = (Id + \frac{\alpha . \Delta t}{(\Delta x)^2} A)U^n
  \tag{$D_{\text{exp}}$}
\end{equation}

et 

\begin{equation}
\label{discretisation:implicite}
  (Id - \frac{\alpha . \Delta t}{(\Delta x)^2} A)U^{n+1} = U^n
  \tag{$D_{\text{imp}}$}
\end{equation}

Dans ces discrétisations, $\text{Id}$ désigne la matrice identité,
$\alpha$ le coefficient de diffusivité thermique, $\Delta t$ le pas de temps,
et $\Delta x$ le pas d'espace

A est la matrice tridiagonale de forme :

$$
A = 
\begin{pmatrix}
-2 & 1  & 0  & \cdots & 0 \\
1  & -2 & 1  & \ddots & \vdots \\
0  & 1  & -2 & \ddots & 0 \\
\vdots & \ddots & \ddots & \ddots & 1 \\
0  & \cdots & 0 & 1 & -2
\end{pmatrix}
$$

$U^n$ est le vecteur solution au temps n :

$$
U^n =
\begin{pmatrix}
u(a,n) \\
u(1,n) \\
u(2,n) \\
\vdots \\
u(b,n)
\end{pmatrix}
\quad \text{et} \quad
U^0 = 
\begin{pmatrix}
u(a,0) = u_{\text{init}}(a) \\
u_{\text{init}}(1) \\
u_{\text{init}}(2) \\
\vdots \\
u_{\text{init}}(b)
\end{pmatrix}
$$

Précisons qu'ici, les exposants sur $U$ indiquent l'indice du temps, et pas une puissance.
Par conséquent, $U^n$ représente le vecteur solution au temps n.
Aussi, les solutions sont des approximations.

Par le cours de Calcul scientifique 1, on apprend aussi la condition de stabilité pour l'équation explicite
est $0 \le \Delta t \le \frac{(\Delta x)^2}{2 \alpha}$

\subsection{Modélisation}

La modélisation métal est la suivante :

\begin{equation}
\label{eq:chaleur-metal}
\begin{cases}
  \frac{\partial u}{\partial t} (x,t) = \alpha \frac{\partial^2 u}{\partial x^2} (x,t), & x \in [a,b],\; t \in [0,T] \\
  u(x,0) = 21, & x \in [a,b] \\
  u(a,t) = 80, & t \in [0,T] \\
  u(b,t) = 21, & t \in [0,T]
\end{cases}
\end{equation}

Et pour la modélisation bois :

\begin{equation}
\label{eq:chaleur-bois}
\begin{cases}
  \frac{\partial u}{\partial t} (x,t) = \alpha \frac{\partial^2 u}{\partial x^2} (x,t), & x \in [a,b],\; t \in [0,T] \\
  u(x,0) = 21, & x \in [a,b] \\
  u(a,t) = 80, & t \in [0,T] \\
  u(b,t) = 21, & t \in [0,T]
\end{cases}
\end{equation}

\section{Implémentation Python}
\label{sec:implementation_python}
% Discrétisation, schémas explicite/implicite, stockage CSR, solveurs (LU, Jacobi, Gauss-Seidel).

\section{Résultats numériques}
\label{sec:resultats_numeriques}
% Profils de température, comparaison des performances des solveurs, analyse de sensibilité.

\section{Reproductibilité}
% Commandes de reproduction, endroit où retrouver le code source, dépendances Python

\section{Déclaration d'assistance par IA}
Dans le cadre de ce devoir, je déclare avoir eu recours à une assistance par IA
pour reformuler certains passages dans un registre plus formel.

J'ai également utilisé ces outils pour identifier des sources documentaires
sur la diffusivité des matériaux, afin d'obtenir une modélisation
aussi fidèle que possible au comportement réel.

\section{Conclusion}
\label{sec:conclusion}
% Bilan, limites et perspectives.

\section{Bibliographie}
\label{sec:bibliographie}

\begin{itemize}
  \item Page Wikipédia de l'acier inoxydable : \href{https://fr.wikipedia.org/wiki/Acier_inoxydable}{wikipedia.org}
  \item Propriétés physiques du bois : \href{https://materiautheque.fr/materiaux/proprietes-du-bois/}{materiautheque.fr}
  \item Conductivité thermique $\lambda$ du bois de bouleau : \href{https://www.engineeringtoolbox.com/thermal-conductivity-d_429.html}{engineeringtoolbox.com}
  \item Page Wikipédia de la capacité thermique massique $c$ : \href{https://fr.wikipedia.org/wiki/Capacit%C3%A9_thermique_massique}{wikipedia.org}
  \item Page Wikipédia de la diffusivité thermique : \href{https://fr.wikipedia.org/wiki/Diffusivit%C3%A9_thermique}{wikipedia.org}
  \item Cours de Calcul scientifique 1 de M1S1 CSMI (Victor Michel-Dansac)
\end{itemize}

\end{document}